\documentclass[10pt,norsk,a4paper]{article}
\usepackage[utf8]{inputenc}
\usepackage[T1]{fontenc}
\usepackage[norsk]{babel}
\usepackage[cm]{fullpage}
\usepackage{color}
\usepackage{parskip,textcomp,amssymb,graphicx}
\usepackage{pdfpages}
\usepackage[stable]{footmisc}
\usepackage{multicol}


\title{Generalforsamling \\ Høsten 2017\\[3cm]
	\includegraphics[width=0.5\textwidth]{cyblogoa3.pdf}\\[-.5cm]}
\date{15.\ november 2017}
\author{Cybernetisk Selskab}

% Blank header, samt footer med side x av y
\usepackage{fancyhdr}
\pagestyle{fancy}
\renewcommand{\headrulewidth}{0pt}
\fancyhead{}
\cfoot{Side~\thepage\ av~\pageref{lastpage}}



\begin{document}

\maketitle{}
\newpage
\tableofcontents{}


\section{Valg av møteleder}
 Jan Furulund blir valg til møteleder ved akklamasjon.

\section{Valg av referent}
 Maximiliano Horta blir valg til referent ved akklamasjon.

\section{Valg av protokollunderskrivere}
 Trine Høyås og Nicolas Papaioannou blir valg ved akklamasjon.

\section{Valg av tellekorps}
 Karl Hole Totland og Jon Even Thorjussen blir valg ved akklamasjon.

\section{Godkjenning av innkalling}
 Flyers til neste år. Godkjent ved akklamasjon.

\section{Godkjenning av dagsorden}
 Godkjent ved akklamasjon.

\section{Semesterberetninger}
\subsection{Semesterberetning ved leder}
\textit{\small Korrigert etter utsendt dagsorden.}\\
\textit{Det blir kort ikke nevnt noe nytt utover det som er lagt ved
skriftlig.}

\textit{Kjære interne, medlemmer og studenter,}
\begin{multicols}{2}

	Dette har vært et bra år for Cybernetisk Selskab. Vi utvidet fadderuka i år, noe som ble en braksuksess. Vi arrangerte storfester de tre siste dagene av fadderuka, og vi hadde fortsatt ikke kapasiteten vi trengte. Selv om det er kjipt med kø, er det veldig postitivt at vi klarer å tiltrekke oss så mange studenter.

	De to andre store arrangementene vi har holdt dette semesteret---dagen@ifi og Ifi-galla---gikk begge uten store problemer. Spesielt dagen@ifi virker det som om at vi har funnet en fin formel på hvordan vi skal arrangere dette sammen med dagen-styret.

	Vi har også dette semesteret startet med en ny type arrangement: \textit{Torsdagsklubben / spritkveld}. Dette er et konsept som har slått bra an. Det er ikke fult hus på disse kveldene, men det kommer mer enn nok folk til at dette er et konsept vi burde fortsette med.

	Personlig har jeg jobbet en del med det pågående prosjektet med å anskaffe nytt kassaapparat. Planen er per dags dato at denne vil bli installert første eller andre uken i desember, og at det blir tatt i bruk etter nyttår. Dette vil medføre enklere økonomirutiner, mindre teknisk kompleksitet, og at de interne i baren får noe nytt å leke med.

	Medlemsmassen vår er på ca.\ det samme som det var i fjor, med 714 semestermedlemmer, noe som vi kan si oss fornøyde med når vi tenker på at vi satte rekord i fjor.

	Helt til slutt vil jeg nevne at Cybernetisk Selskab fyller 50 år i 2019. Jeg har sammen med noen pamper og andre interesserte ifi-studenter startet i det små for å finne ut hva vi ønsker å få til. Vi ønsker å danne et ``CYB-50''-utvalg som skal ta seg av organiseringen av jubileumet, og at det skal bestå av både gamle og nye studenter. Jeg håper at dette er formelt på plass like etter nyttår. Vi har også snakket med Fordelingsutvalget, som er positive til å støtte prosjektet så lenge vi skriver en intensjonsmelding på hva vi ønsker å gjennomføre, samt sørger for at dette er et arrangement for hele instituttet.

	Vi sees over nyåret!

	\textbf{Andreas Hansen},\\
	leder \\
	\date{10.\ november 2017}
\end{multicols}

\newpage


\subsection{Semesterberetning ved kjellermogul}
\textit{\small Korrigert etter utsendt dagsorden.}\\

\textit{Det blir kort ikke nevnt noe nytt utover det som er lagt ved
skriftlig}\\
\begin{multicols}{2}
	Da nærmer atter et semester seg sin slutt, og ikke annet kan egentlig sies enn dette har vært et svært godt år for CYB og Escape.

	I fadderuka i år kjørte vi utvidet område ikke bare på fredagen, men også på onsdagen og torsdagen. Selv om dette hadde sine utfordringer, gikk det bra for å være et pilotprosjekt, og vi endte opp med en markant høyere omsetning for hele uka, samt at mange flere funksjonærer stilte opp enn året før. Årets fadderstyre fokuserte også veldig på å opprettholde en god kommunikasjon med oss, og hjalp til med å sørge for at alle faddere hadde arbeidsoppgaver, der hvor å jobbe i bar eller som vertskap gikk innunder dette.

	Vi har også hatt enda flere arrangementer enn før, hvor nesten alle ukene i semesteret har hatt noe som har foregått. Dette er egentlig noe som bare blir bedre med hvert år, som kanskje er selvsagt. Dagen@ifi gikk også bedre enn noen gang, hvor vi i år økte omsetning på ca. 50 000,- i forhold til i fjor. Det er tydelig at vi med tiden har begynt å få disse arrangementene mer og mer under huden, og har nå ganske mye informasjon om hvordan man faktisk arrangerer de større arrangementene på instituttet. Jeg kan også meddele at renhold på en generell basis er veldig fornøyde med oss når det gjelder så og si alle de store festene vi arrangerer, noe som jo er positivt.

	Torsdagsklubben er også et nytt prosjekt der vi prøver å jevnt ha drinkkvelder i løpet av semesteret, både for å trene opp folk i hvordan å jobbe med drinker (til større arrangementer hvor vi trenger flinke folk) og for å gi et tilbud til de som er glad i drinker framfor annen alkoholholdig væske for å gi rom for et enda større publikum. Selv om---som også Andreas påpeker---at oppmøtet ikke nødvendigvis er så stort, så er det noe jeg tror bare kommer til å bli bedre med tiden hvis vi klarer å fortsette med slike arrangementer. I samme slengen kan det også nevnes at vi har hatt såpass mange arrangementer at det har ført til at det er mange utlånsforespørsler som vi har måttet takke nei til, fordi det rett og slett ikke har vært plass på ukeplanen. På mange måter vil jeg si at det bare er en positiv indikasjon på Escape sin rolle som møte- og knutepunkt.

	Web- og driftgruppene ble slått sammen til X-gruppa, som jeg tror viser at det er---og kommer til å bli---et enda større fokus rundt både det å få inn flere interne i gruppa, og jobb med utvikling av internsystemet, wikien og nettsidene våre.

	For å følge opp det jeg nevnte sist semester om kjelleravtalen, har Cosa Nostra generelt jobbet ganske lite med den nye (foreløpige), utover dokumentet vi sendte til UiO med krav om ønskede forandringer i forrige kjelleravtale. Så vidt jeg vet sikter universitetet sentralt fremdeles på å fullføre den nye avtalen innen første halvdel av 2018.

	Til slutt vil jeg egentlig bare si takk for meg, og takk for et supert år med to autonome og produktive kjellerstyrer, samt en fantastisk internbase som er de som faktisk får ting til å gå rundt. Jeg er---som sikkert er sagt av mange før meg---veldig positiv til hva CYB kan få til framover, og har fått til på ganske få år, og håper at vi en dag klarer å bli ikke bare en av de største, men den største kjellerforeningen ved Blindern.

	Vi sees i Escape!

	\textbf{Karl Hole Totland}, \\
	kjellermogul \\
	\date{14.\ november 2017}
\end{multicols}

%\newpage

\section{Kasserer orienterer}
\subsection{Regnskapets tilstand}
Kasserer orienterer om regnskapet.

Ligger godt an med regnskapet. Har utestående 61 tusen og alle påmeldinger har
ikke blitt registrert fra skitur og penger fra FUI. Vi har derfor minus og etter
det er ført inn så kommer vi til å ligge med 20~30K i pluss. Gitt at det ikke
kommer noe uventede utgifter ligger vi godt an med budsjettet. Noen punkter
har gått over på utgifter men samtidig gått over på inntekter og gjort seg
bedre enn motsatt.

Det eneste som har gått over budsjett uten at det har regulert seg selv er
IFI-Galla der det ligger an til et tap med 20~25tusen. Det mistenkes at mindre
støtte til gallaen som fører til tapet.

\textit{Det ligger ca 22 flasker som ikke er drukket opp, dette er ikke
flyttet over regnskapsmessig og vil minske dette tapet.}

Uspesifisert -38000 i regnskapet skyldes ny kaffekvern som ikke har fått egen
konto.


\subsection{Revidert budsjett for høsten 2017}
Kasserer orienterer om budsjett.

Budsjett for neste semester er det gjort noen innskrenkinger på blant annet
Magic Gathering da det er lite utgifter der. 
En ny post blir det nye kassesystemer der det vil komme en utgift på ca 10K i
året.
Varer har fått økning i budsjett og det er budsjettert med å gå med 11k i
pluss i 2018.


\section{Kontigentfastsettelse}
Hovedstyret foreslår å endre medlemskontigenten fra kr. 30,- til kr. 40,-,
fra og med neste semester. \\
Hovedstyret orienterer rundt motivasjonen bak dette.

Den har ikke blitt rørt på mange år. Kassesystem og Cyb-50 er utgifter som vi
ser vil øke, og istedenfor å øke vareprisen på ting vi selger.. Escape er en
av de med lavest utgifter blant foreningene. De som snur i døren pga inngang
på 40 kr vil også snu for kr 30. 


\textbf{Det er vedtatt ved flertall}

\section{Valg}

\subsection{Hovedstyret}
Man velges inn i hovedstyret for ett år av gangen.

\begin{multicols}{2}
\subsubsection{Kjellermogul}
\textit{Kjellermogul Karl gir en kort orientering om vervet}

Trine, tidligere kaffesjef og føler at kjellermogul er et naturlig vis
videre da hun ikke føler seg ferdig med verv. Hun har god oversikt over ting
som må gjøres og ting som må gjøres rundt i kjelleren.

Adrian Har sittet som internansvarlig i 1,5 år og ønsker å bidra på en
annen måte en det han har gjort. Ønsker å få på plass nye rutiner av nye
funker og bedre oppfølging av kjellerstyret. Adrian ønsker å ha mer ansvar ved
større arrangementer og bidra mer.

\textit{Tid til vervet}\\
T: skal ta opp fag og har god tid\\
A: Har fått tid til ting tidligere og regner med det samme gjelder videre

\textit{Stressmestring}\\
A: Er relativ rolig og får ting gjort ordentlig, ikke redd for å spørre om
hjelp hvis han ser det ikke er tid til å gjøre det selv\\
T: Noe hun må jobbe mer med under krise, men har til vanlig god kontroll

\textit{Delegere arbeid - Hvor flink}\\
A: Når han selv ser at han ikke har kapasiteten til å fikse ting selv, så er
han selv flink til dette, og flink til å spørre hvordan det går i ettertid og
følge opp\\
T: Føler hun er flink til dette og har blitt flinkere til dette i løpet av
året som har fått.

\textbf{Adrian Helle vinner valget og blir valgt.}

\subsubsection{Kasserer}
\textit{Kasserer Stian gir en kort orientering om vervet}\\
Stian stiller til gjenvalg\\
\textbf{Stian Rustad velges ved akklamasjon.}
\subsubsection{Arrangementssjef}
\textit{Arrangementssjef Elise gir en kort orientering om vervet}\\
Elise stiller til gjenvalg\\
\textbf{Elise Lystad velges ved akklamasjon.}
\subsubsection{Promoteringssjef}
\textit{Promoteringssjef Jon Even gir en kort orientering om vervet}\\
Kathrine stiller som Promoteringssjef. Kathrine er førsteårsstudent
som tidligere har studert medier og kommunikasjon og har designerfaring derfra.
Har også tidligere foreningserfaring med design. \\
\textbf{Kathrine Hoel velges ved akklamasjon.}
\newpage

\subsubsection{Internansvarlig}
\textit{Internansvarlig Adrian gir en kort orientering om vervet.}\\
Lilly har vært barsjef og tror hun kan bidra til internmiljøet. Hun
ønsker at Escape skal være et hyggelig sted å være og har selv hatt denne
erfaringen og ønsker at de interne skal ha samme følelse etter å blitt
interne.\\
Lee er ofte i kjelleren og ønsker å sørge for at Escape er et sted
du kan komme til og føle deg velkommen. Hun ønsker å være en synlig resurs
person som man kan komme til uavhengig av verv man måtte ha eller bare er
funk. Hun er en person som stiller opp ved behov.\\
Henrik Økonomifunk og tidligere erfaring fra f.eks Navet og SVFF. Ønsker å
involvere en større del av internmassen på interne arrangementer.

\textit{Spesefikke ting de øsnker å få gjort}\\
Lee: Ønsker å gjøre det lettere å finne ut hvem man kan/skal ta kontakt med
sånn at de har en kontaktperson. Det tar for lang tid å finne den
rettepersonen og hun øsnker å være bindeledd
Lilly: Ønsker å få kosetirsdag til et mere sosialtarrangement hvor flere kommer
og forslagspostkasse
Henrik: Prøve å få flere til kosetirsdag og få flere til å være interne utover
det de gjør for å være sosiale på tvers av personer

\textit{HR-saker med en ukjent person mot en person du kjenner godt}\\
Lee: Prøver å være objektiv som mulig, hun vet det er vanskelig men det er
nøvendig og være det da internansvarlig har dette ansvaret og vil derfor være
objektiv. Hun ønsker og få oversikt over intensjonen og hensiktet fra begge
parter for å vite hvorfor dette skjedde og resultatet for å kunne megle best
mulig

Lilly: Må kunne kommunisere med begge parter, og ser på det at nærevenner eller
uansett hvem som dette gjelder så skal hun greie å stille seg som en
meglerpart og åpne for kommunikasjon

Henrik: Ser på dette som noe han skal greie uten problem og kommer ikke til å
forskjellsbehandle


\textit{Det er ikke bare som sørger for internmiljøet blir bra, hvordan skal
dere ta det opp med mellomledere?}\\
Lee: Ønsker regelmessig oppfølging av de som jobber skal bli hørt, og være en
person som kan kommunisere deres behov.
Henrik: Spørre de interne om de føler de får nok informasjon og snakke med
mellomledere/styremedlemmene og følge disse opp. Ikke kun fikse nedenfra
oppover men begge veier.
Lilly: Sørge for at de i baren som jobber lite å få dem introdusert til
internmiljøet ved hjelp av kosetirsdager og rett og slett hilse på dem

\textbf{Lee Kvåle blir valg til internansvarlig}

\end{multicols}
\newpage
\subsection{Kjellerstyret}
Alle verv--\textit{utenom økonomiansvarlig}--som er til valg i \\kjellerstyret gjelder for ett semester av gangen.

\begin{multicols}{2}
\subsubsection{Økonomiansvarlig}
\textit{Økonomiansvarlig Svein orienterer kort om vervet}\\
Maela er førsteårsstudent og ønsker å stille. Jobber i kaffen og har erfaring
derfra og vet hvordan priser og slikt er. \\
\textbf{Maela Le Borgne blir valgt ved akklamasjon}
\subsubsection{Innkjøpsansvarlig}
\textit{Innkjøpsansvarlig Tor-Aksel orienterer kort om vervet}\\
Tor-Aksel stiller til gjenvalg. Skal også ta med seg en innkjøpsfunk som kan
lære bort det han har lært gjennom vervet og jobb.\\
\textbf{Tor-Aksel Solberg blir valgt ved akklamasjon}
\subsubsection{Barsjef}
\textit{Barsjef Lilly orienterer kort om vervet}\\
Lilly stiller til gjenvalg. Ønsker å fortsette at baren skal være et
sted man ønsker å være og fortsette å drifte baren slik det har blitt gjort.
Marie har jobbet i baren siden V16, både i bar og kafé. Ønsker å
sette fokus på opplæring av funker og skjenkemestere sånn at de får en følelse
at de bidrar med noe og gjør en god jobb samtidig.\\
\textit{Skal dere holde Innkjøpsansvarlig litt i ørene}\\
L: Det er ikke plass til mere
M: Tenker det er viktig å ha god kommunikasjon med Innkjøpsansvarlig sånn at
vi alltid har ting som vi skal ha. Både av sprit og øl når det trengs.

\textit{Hvilke konkrete tiltak for oppklæring av SM}\\
L: Samtaler med SM med den som ga og fikk opplæring og sørge for at de som gir
opplæring har god erfaring før de gir opplæring
M: Tenker det er en god idee, men samtidig se hvem av de i baren som kan greie
å ta på seg et slikt ansvar og passer til vervet, ikke gi dette til hvem som
helst.\\
\textbf{Marie Reichelt velges ved flertall}

\subsubsection{Kafésjef}
\textit{Kafésjef Trine orienterer kort om vervet}\\
\textbf{Håkon} er tidligere kaféfunk og trives veldig godt der. Tror han
kommer til å gjøre en god jobb. Han er en person som liker å ta på seg mye
ansvar\\
\textbf{Amanada} ønsker å fortsette det fantastiske arbeidet Trine har gjort og er
alltid her og har derfor gode muligheter til å gi en ordentlig oppklæring av
folk i kafeen. Ønsker også ha god kommunikasjon med Innkjøpsansvarlig.\\
\textit{Noe spesielt å gjennomføre}\\
H: Fornøyd slik det er nå og har store sko å fylle, men har ikke noe
spesielt.\\
A: Ting fungerer veldig bra og fortsette med det som gjøres nå

\textit{Mer deppressiv musikk i kafeen?}\\
A: Kan spille av mere hvis dere øsnker det, og mindre\\
H: Kan gjøre det dere vil

\textit{Droppe forelesnigner hvis dere må?}\\
H: ja\\
A: Ja

\textit{Forhold til delegering og det å finne folk til å være reserve og
funksjonærer for å hindre at dere selv brenner dere ut}\\
A: Ha god dialog og en god reserveløsning
H: Kommer absolutt til å delegere der det er mulig, men er villig til å ta på
seg mye hvis det er behov, samt opprettholde dialogen

\textit{latte-art}\\
H: Er egentlig ikke så veldig flink, men gjør sitt beste
A: Lager gjerne penislatte hvis dere ønsker det.\\
\textbf{Håkon Bliksås Carlsen velges til kafésjef ved flertall}

\subsubsection{Teknisk ansvarlig}
\textit{Tekniskansvarlig Nicolas orienterer kort om vervet}
Sindre stiller da han har drevet med lignende ting i 7 år.\\
\textit{Tekniskansvarlig til faktisk dokumentasjon}
Nicolas skal være her ca 2 år til, har tid til å få det dokumentert\\
\textbf{Sindre Sjøstrand Brurberg velges ved akklamasjon}

\subsubsection{DJ-sjef}
\textit{DJ-sjef Vincent orienterer kort om vervet}
Nicolas stiller til DJ sjef, har tidligere vært teknisk sjef i 2 år og vært
aktiv DJ siden han startet. \\
\textbf{Nicolas Harlem Eide velges ved akklamasjon}

\subsubsection{Utlånsansvarlig}
\textit{Utlånsansvarlig Benjamin orienterer kort om vervet}
Benjamin stiller til gjenvalg\\
\textit{Hvorfor gjenvalg}
Konfortabelt og god kommunikasjon med SM, også mulig til å sette seg selv på
jobb.
\textit{Fortsetter du i SPF}
Kommer ikke til å ta ledervervet der, men fortsetter 
\textbf{Benjamin Breiby velges ved akklamsjon}

\newpage
\subsubsection{Arrangementskoordinator\footnotemark}
\textit{Kjellermogul Karl orienterer om vervet}
Skal avlaste Arrangementsstyret og kjellerstyret med å koordinere
arrangementene. Skal ikke ta over noe ansvaret men hjelpe til for
arrangementerer som skal innebære bar.

\textit{Edgar} stiller til vervet. Han ble informert om vervet og lenge vært i Cyb. Vært en pådriver for
Torsdagklubbene og alle store arrangementer som Cybernetisk Selskab har organisert. Har
aldri hatt et ordentlig verv, men alltid han har alltid bidratt og tatt i
ekstra på både i baren og organisering av festene.

\textit{Shakere}
Boston

\textit{Tidligere styreverv}
Paintballforeningen, der han støttet

\textit{Hva er det du som dette styrevervet}
Ble mye å gjøre for øvrige styremedlemmer, så det blir bare en avlastning.\\

\textit{Storcashturerer 1 gang i mnd blir for mye når man ikke har noen som vet
hvordan spritkvelder går osv. Dette vervet blir lettere å skalere opp
arrangementer som er litt større enn fredagspub. Dette har pleid å være noe
som har bare blitt kastet rundt uten at noen har egentlig tatt ansvaret.}\\

\textit{Hva annet en spritkvelder er det du skal gjøre}
Håper om å gjøre ting mer fleksibelt. Noe man gjør etter hva man ønsker. Han
ønsker at dette skal være avlastning sånn at styrevervene har bare blitt
større etter som cyb har blitt større med årene.
Har ingen konkrete ting han ønsker å gjennomføre, men ønsker å gjøre noe

\textbf{Edgar Andorsen blir valgt ved akklamasjon}

\end{multicols}
\footnotetext{Tentativ vervtittel.}

\newpage

\section{Vedtektsendringer}
Følgende forslag til vedtektsendringer ligger fremme.

\subsection{Endring i forbindelse med krav om Ifi-foreninger}

For å være en offisiell Ifi-forening har Ifi satt sammen noen vedtekter som studentforeningene må følge.
Per dags dato følger ikke Cybernetisk Selskab §1c\footnotemark. Det ble lagt frem et forslag av hovedstyret på vårens generalforsamling om å endre våre vedtekter, slik at vi ikke bryter med Ifi sine krav.
Dette ble da ikke vedtatt da forsamlingen konkludere med at de nye vedtektene var for åpne og uklare. Forsamlingen ønsket at hovedstyret skulle komme med et nytt forslag til neste generalforsamling, som blir denne, høsten 2017.

\footnotetext{Etter manglende vedtektsfesting av to-tredelers Ifi-studenter i hovedstyret.}

Om vedtektene ikke blir vedtatt vil hovedstyret starte søknadsprosessen for å få et unntak fra §1c av \textit{Vedtekter for Ifi-forening}, slik som §3c åpner for. Begge endringsforslagene hører sammen og stemmes over i lag.

Hovedstyret er innstilt på å starte søknadsprosessen etter §3c.

\subsubsection{Forslag til ny paragraf: $§5h$}
\begin{quote}
	\begin{enumerate}
		\item[§5h] Minimum to-tredeler av hovedstyret skal bestå av studenter som enten (1) er tilknyttet et eller flere studieprogrammer ved \textit{Institutt for informatikk}, eller (2) som tar ett eller flere IN- eller INF-fag.
	\end{enumerate}
\end{quote}

\subsubsection{Forslag til endring av paragraf: §8a}
\begin{quote}
	\begin{enumerate}
		\item[§8a]
			Valgbare til alle styrer er alle medlemmer i henhold til §2 og informatikkstudenter ved Universitetet i Oslo. \textcolor{red}{Om valg av kandidat bryter med §5h, blir valget gjort ugyldig og det blir gjenvalg. Om ingen andre kandidater stiller blir vervet etterfylt av hovedstyret i henhold til §5f.} Alle valg avgjøres \textcolor{red}{på generalforsamling} ved simpelt flertall.
	\end{enumerate}
\end{quote}

Tidligere presentert og ble nedstemt forrige Generalforsamling.
Ble nedstemt grunnet at §5h var for løst og hva gjorde man hvis man hadde et
verv som allerede var valg og vanskeliggjorde neste valg da det var en viss
mulighet for at noen ikke kunne stille til neste valg siden det allerede var
valgt folk som ikke var ifistudenter. Dette skulle løses ved §8a.

\textit{Har det vært HS der dette ikke oppfylles?} Så langt bakover har dette aldri
vært et problem tidligere\\
\textbf{Saksopplysning:}\textit{Etterfylling er ikke et valg, og står forklart
i lovene.}

pressiseringen i §8a er gjort for å få setningen ikke til å være merkelige og pga
plassering av vedtekter som gjør at det kan missforstås uten dette.

\textit{Kan Instituttet tvinge oss til å vedta dette?}
De kan ikke dette, de kan si vi ikke er en ifi-forening, ikke tilgang til FK,
ikke tilgang til bygget og kommer ikke på hjemmesiden. Dette er kontraktfestet
med UiO gjennom kjellerkontrakten.

\textit{Mest sannsynlig kommer de ikke til å gjøre dette pga det dårlige lyset ifi blir satt
i av å gjøre dette.}

\textbf{Hvis vi ikke er ifi forening så er vi fortsatt Institutt foreningen
ved ifi}

\textit{De andre ifi foreningene har spesifisert at styret må være 2/3 ifistudenter}


\textbf{paragrafene godkjennes ikke.}

\section{Æresmedlemskap}
Ingen utnevnes til æresmedlemskap

\section{Utdeling av pins}\label{lastpage}
Internansvarlig deler ut pins på vegne av arkivar.

De som får pinnen er følgende:\\
\textbf{Jon Even}\\
\textbf{Elise}\\
\textbf{Trine}\\
\textbf{Svein}\\
\textbf{Vincent}\\
\textbf{Tor-Aksel}\\
\end{document}
